\scanpage{14}
% archive foliation: 13
$k$ anneau (com. unit.)

\uline{Déf 1} Un $k$-module $\frac12$ simplicial $L_*$ est dit
$\partial$-pseudo-cohérent (resp. strictement $\partial$-pseudocohérent,
resp. $\partial$-parfait, resp strictement $\partial$-parfait)
si $L^!_*$ (le complexe de chaînes associé
par Dold-Puppe) est pseudo-cohérent
(resp. str. ps. cohérent, resp parfait, resp.
strict.t parfait).
\margin{\unsure{Donc on a par déf a) ci-dessous}}

\uline{Prop 1} a) $L_*$ est str. $\partial$-pseudocohérent
(resp. str. $\partial$-parfait) si et seulement si les
$(L^!_n)$ sont projectifs de t.f. (resp. et
nuls sauf un nb fini); \deleted{\illegible{3}}

b) $L_*$ est str. $\partial$-ps. coh. ssi les $L_n$
sont projectifs de t.f.

c) $L_*$ est str. $\partial$-parfait ssi $L_*$ est
extension successive \inserted{finie} de Modules de
la forme $\Lambda^i\Psi_*\otimes_{\mathbf Z}M_i$
\inserted{$(i\geq0)$}, où les
$M_i$ sont $k$-modules pr. de t.f.
(l'ordre des facteurs dans une filtration
croissante \unsure{par exemple|précédente} pourra d'ailleurs
être choisi par ordre des $i$
croissant), ou encore \deleted{\unsure{ssi}}
(\unsure{supposant $k$ connexe à spécifier comme une hypothèse}).
\margin{ssi les $L_n$ sont proj. de t.f. et
$f_{L_*}(n)=\operatorname{rg}_kL_n$ \unsure{une fonction
polynomiale pour $n$ grand} \unsure{(on}\illegible{3}\unsure{telles fonctions)}}

d) Si $k$ est noeth., $L_*$ est $\partial$-
pseudo-cohérent (ssi les $L^!_n$ sont
de type fini sur $k$, \unsure{ou encore}) ssi
les $L^{\unsure{\natural|\sharp}}_n$ le sont.

e) Si $k$ est noeth. régulier, $L_*$ est
$\partial$-parfait ssi les $L_n$ sont de t.f.
et \deleted{\unsure{$L_*\otimes K$}} $\forall$ $K$ corps résiduel, la fonction
